\subsection{Persistent-Marker Spherical Direct Simulation Monte Carlo
 (SphDSMC)}
\label{sec:spherical-dsmc}

\noindent\textbf{Maturity and validation.}
The constant-cross-section self-interacting dark matter (SIDM) configuration
implements the persistent-marker, sorting, collision-shell, and event-count
conventions of the finite-time benchmark described below.  Following a halo
through gravothermal core collapse additionally requires convergence in marker
count, timestep, shell population, collision locality, and random seed, with
conservation and restart checks throughout the evolution.  Velocity-dependent
cross sections require their own finite-time and long-time comparisons.

\begin{figure}[htbp]
\centering
\textbf{Physics flow: persistent-marker spherical DSMC module}
\par\smallskip
\begin{tikzpicture}[node distance=0.22cm]
\node[physicsflow] (sd1) {Sample spherical markers from the initial distribution, or restore their individual masses and orbits from a restart};
\node[physicsflow,below=of sd1] (sd2) {Sort the markers by radius and use their radial order to calculate enclosed mass and gravity};
\node[physicsflow,below=of sd2] (sd3) {Choose one physical timestep from the orbital and collision-rate limits};
\node[physicsflow,below=of sd3] (sd4) {Apply a half gravity kick, move every marker, sort the crossed markers again, and apply the second half kick};
\node[physicsflow,below=of sd4] (sd5) {Retry any inaccurate close central passage with the regularized orbit calculation};
\node[physicsflow,below=of sd5] (sd6) {Divide the newly sorted halo into local radial collision shells no wider than the mean free path when resolved};
\node[physicsflow,below=of sd6] (sd7) {Choose trial pairs, accept them at the physical collision rate, and scatter accepted velocities elastically};
\node[physicsflow,below=of sd7] (sd8) {Check mass, energy, orbit accuracy, and collision locality; then write profiles or a restart and continue};
\draw[physicsarrow] (sd1) -- (sd2);
\draw[physicsarrow] (sd2) -- (sd3);
\draw[physicsarrow] (sd3) -- (sd4);
\draw[physicsarrow] (sd4) -- (sd5);
\draw[physicsarrow] (sd5) -- (sd6);
\draw[physicsarrow] (sd6) -- (sd7);
\draw[physicsarrow] (sd7) -- (sd8);
\end{tikzpicture}
\end{figure}

\subsubsection{Purpose and physical picture}

SphDSMC solves the static Boltzmann--Poisson equations (PIC0b)--(PIC0c), with
\(H=0\), the full halo density as the Poisson source, and elastic binary
scattering in \(C[f]\).  It uses the same spherical phase space measure
(PIC0a).  The marker ordering supplies the enclosed mass directly, so the
force calculation implements the spherical Poisson equation without a
separate gravity mesh.  The orbital update and local collision sampling then
advance the two terms in the kinetic equation.

This solver represents the halo by a large set of computational
particles, called \emph{markers}.  Each marker keeps its radius, radial
velocity, angular momentum, and positive mass throughout the run.  A fresh
standard sample is equal mass, but the gravity, shell, collision, diagnostic,
and restart paths accept unequal masses.  Gravity moves the markers along
spherical orbits.  Collisions select nearby whole-object pairs and scatter
their continuous three-dimensional velocities.
An origin crossing is always reflected as
\(r\mapsto|r|\), \(v_r\mapsto-v_r\).  There is no outer wall: \code{r_max}
is the outer edge of the reported profile bins, not a reflecting boundary,
and a marker outside it remains part of the gravitational state.

It follows the SphDSMC construction of
\href{https://doi.org/10.1093/mnras/stag2015}{van den Bosch and Tokayer (2026)}.
The input \(f(r,w,L)\) supplies the probability distribution from which the
initial markers are drawn.  After this sampling step, the markers themselves
are the dynamical state.  The task selector and runnable interface are given
only after the state, gravity, timestep, and collision laws below.

One timestep proceeds as follows:
\begin{enumerate}
\item Sort the markers by radius and compute the spherical gravitational
      force from their radial order.
\item Apply half of the gravitational velocity kick.
\item Move every marker to its new radius.
\item Sort again, because markers may have crossed one another.
\item Apply the second half of the gravitational kick using the new radial
      order.
\item Divide the sorted markers into radial collision shells.
\item Draw the physical number of collision trials in each shell, accept
      pairs with the Boltzmann probability, and scatter every accepted pair.
\item Write profiles, conservation budgets, and restart data at their
      requested intervals.
\end{enumerate}

The orbit and output logic lives in
\path{cpu/include/ev/spherical_dsmc_solver.hpp}.  The pair-selection and
finite-angle scattering logic lives in
\path{cpu/include/ev/spherical_dsmc_collision.hpp}.

\subsubsection{How SphDSMC differs physically and numerically from WPIC}

The two modules approximate the same collisionless Vlasov--Poisson equation
between encounters, but their Vlasov updates are not identical numerical
algorithms.  In WPIC, unequal-weight markers deposit their mass onto a fixed
radial mesh; the spherical Poisson equation is solved on that mesh, and the
interpolated mesh force moves each marker.  In SphDSMC, the persistent markers
are sorted by their instantaneous radius, the actual marker masses are
prefix-summed in that order to obtain \(M(<r)\), and each marker receives the
corresponding spherical force directly.  SphDSMC therefore needs a stable
radius sort after every drift.  WPIC needs no radial sort for its gravity
solve, although its material-collision path may sort markers when constructing
adaptive collision shells.  Both use a time-centred kick--drift--kick orbit
update, so they solve the same physics while having different force noise,
spatial smoothing, and convergence errors.

Their elastic SIDM collision cores share the same physical ingredients:
local no-time-counter pair trials, the relative speed \(u\), a possibly
velocity-dependent pair area \(\Sigma_{ij}(u)\), and exact unequal-mass
centre-of-mass scattering.  ``No-time-counter'' means that trial pairs are
drawn from a safe rate bound rather than by advancing a separate clock for
every possible pair.  The complete SphDSMC pair remains two markers and the
present SphDSMC module has no inelastic capture channel.  Its collision
timestep is part of the one global physical timestep and elastic scattering
is applied once per accepted global step.

WPIC additionally has the Plummer-subhalo encounter operator.  It can either
scatter two complete unequal-mass subhalos or capture them into one remnant of
mass \(m_i+m_j\), placing the lost relative orbital kinetic energy into the
remnant's internal energy and velocity dispersion.  WPIC also separates the
ordinary Vlasov CFL timestep from the outer collision-operator deadline and
from the collision sampler's internal probability-limiting substeps.  Thus
the SphDSMC and WPIC collision laws can use the same elastic SIDM cross section,
but the calling cadence, collision-region construction, gravitational force
estimator, and available inelastic physics are not the same.

\subsubsection{Marker mass and initial sampling}

One marker stores
\codesource{SphDSMC marker state}
{\path{cpu/include/ev/spherical_dsmc_solver.hpp}}
{Defines the persistent radius, radial velocity, angular momentum, identifier,
 and represented mass carried by each spherical marker}
\begin{lstlisting}[language=C++,basicstyle=\ttfamily\scriptsize]
struct SphericalDSMCParticle {
    double r;
    double vr;
    double L;
    std::uint64_t id;
    double mass;
};
\end{lstlisting}
\code{SphericalDSMCParticle} is the C++ data record for one persistent
marker: it holds the marker's current phase space coordinates, immutable
integer identifier, and represented mass.  Here \(v_r\equiv w\), and
\[
 L=r v_t,\qquad v_t^2=v_\theta^2+v_\phi^2.
\]
The standard initializer gives every marker the same mass
\[
 m_p=\frac{M_{\rm represented}}{N_p}.
\]
Normally \(M_{\rm represented}\) is the integral of the carrier grid.  If the
source particle model supplies an independently known total mass,
\code{bp_dsmc_total_mass} sets that mass explicitly.  It changes only \(m_p\):
the normalized cumulative mass distribution, and therefore the sampled
phase space shape, stays fixed.

For a spherical distribution \(f(r,v_r,L)\), integration over position and
velocity angles gives
\[
 \dd M=8\pi^2 Lf(r,v_r,L)\,\dd r\,\dd v_r\,\dd L.
\]
Consequently, the piecewise-constant mass represented by one input cell is
\[
 \Delta M_{ijk}
 =8\pi^2L_k f_{ijk}\,
   \Delta r_i\,\Delta v_{r,j}\,\Delta L_k.                 \tag{S1}
\]
In equation (S1), \(i\), \(j\), and \(k\) index the radial,
radial-velocity, and angular momentum cells, respectively.  Elsewhere \(p\)
denotes an arbitrary marker, \(q=0,\ldots,N_p-1\) denotes its position in the
stratified initial sample, \(i\) denotes zero-based radial rank after sorting,
and \(s\) denotes a collision shell; each local use is stated explicitly.
The radial control-volume array has its innermost face at exactly \(r=0\).
The velocity and angular momentum widths are the same trapezoidal
control-volume widths used by the grid solver.

The code forms a cumulative table of the positive cell masses in (S1).
It then places \(N_p\) evenly spaced target masses along that table, using one
random offset \(\xi\in(0,1)\).  Marker \(q\) selects the cell containing
\[
 M_q=M_{\rm represented}\frac{q+\xi}{N_p},
 \qquad q=0,\ldots,N_p-1.
\]
This gives exactly \(N_p\) initially equal-mass markers and spreads the
selected cells evenly across the cumulative mass distribution.  The mass is
then stored on each marker.  Restart versions 4--7 preserve it individually;
collision code never infers equality from the initializer.

Within the selected cell, \(r\) and \(v_r\) are uniform.  Since the measure in
\(L\) is \(L\,\dd L=\tfrac12\dd(L^2)\), the sampler is uniform in \(L^2\):
\codesource{SphDSMC grid-to-marker sampler}
{\path{cpu/include/ev/spherical_dsmc_solver.hpp}}
{Samples continuous spherical phase coordinates from a selected positive-mass
 carrier cell, including the correct angular momentum measure}
\begin{lstlisting}[language=C++]
particle.r = r0 + (r1 - r0) * rng.uniform_open();
particle.vr = w0 + (w1 - w0) * rng.uniform_open();
particle.L = std::sqrt(
    L_lo * L_lo
    + (L_hi * L_hi - L_lo * L_lo) * rng.uniform_open());
\end{lstlisting}
The transformation from \(\dd^3x\,\dd^3v\) to
\(\dd r\,\dd v_r\,L\,\dd L\) cancels the radial \(r^2\) factor, yielding
uniform sampling in \(r\) inside the selected cell.

\subsubsection{Radius ordering, spherical force, and total energy}

The marker array is maintained in nondecreasing radius order.  For zero-based
radial rank \(i\), let \(m_i\) be that marker's mass, and let \(j<i\) denote
all strictly interior ranks.  The fixed central point mass is
\(M_\bullet=\code{central_mass}\); it is zero when no central object is
present.  The shell-centered enclosed force mass is
\[
 M_i^{\rm enc}=M_\bullet+\sum_{j<i}m_j+\frac12m_i.         \tag{S2}
\]
The characteristic equations are
\[
 \dot r_i=v_{r,i},\qquad
 \dot v_{r,i}=\frac{L_i^2}{r_i^3}
              -\frac{GM_i^{\rm enc}}{r_i^2},
 \qquad \dot L_i=0                                         \tag{S3}
\]
between collisions.

Equation (S2) is the spherical-shell analogue of centering a finite-volume
force in a mass cell.  The corresponding diagnostic kinetic and potential
energies are
\[
 K=\sum_i\frac{m_i}{2}
 \left(v_{r,i}^2+\frac{L_i^2}{r_i^2}\right),               \tag{S4}
\]
\[
 W=-G\sum_i\frac{m_iM_i^{\rm enc}}{r_i}.                  \tag{S5}
\]
The half-self contribution in (S5) is
\(-Gm_i^2/(2r_i)\) for the represented thin spherical shell and is consistent
with the centered force (S2).  The reported Newtonian virial ratio is
\(Q_{\rm vir}=2K/|W|\); values near one indicate global virial balance, but
this single scalar does not prove local equilibrium.

The ordinary update is kick--drift--sort--kick:
\codesource{SphDSMC gravity step}
{\path{cpu/include/ev/spherical_dsmc_solver.hpp}}
{Applies time-centred spherical gravity and rebuilds enclosed mass after marker
 crossings before the endpoint kick}
\begin{lstlisting}[language=C++]
particles[i].vr += 0.5 * dt
    * spherical_dsmc_acceleration(
          particles[i], enclosed_force_mass[i], p);
particles[i].r += dt * particles[i].vr;
spherical_dsmc_sort_particles(particles);
rebuild_enclosed_force_mass_from_sorted_marker_masses();
particles[i].vr += 0.5 * dt
    * spherical_dsmc_acceleration(
          particles[i], enclosed_force_mass[i], p);
\end{lstlisting}
The intermediate sort is physical: after shell crossings, the second kick
must use the new enclosed rank.  A one-thread paper-exact run uses Frank and
Tokayer's literal insertion pass after every evolved step.  A threaded
paper-exact run produces the same stable radial ordering with either that
literal pass or the stable block-and-merge decomposition described below.
Only the one-time initial ordering uses a general stable ordering method.
The generic VENUS path may use insertion, block-and-merge, or radix sorting
according to the measured disorder.

\paragraph{Insertion sort: repairing the order left by a small step.}
Insertion sort walks from the second marker to the end of the array.  At
position \(i\), it temporarily saves the marker as a \emph{key}, shifts every
preceding marker with radius larger than the key one place to the right, and
inserts the key into the empty position.  Equal radii are not shifted past one
another, so the ordering is \emph{stable}: markers with equal sort keys retain
their previous relative order.  In source form the essential operation is
\codesource{SphDSMC stable insertion sorter}
{\path{cpu/include/ev/spherical_dsmc_solver.hpp}}
{Repairs a nearly ordered radial marker range without changing the relative
 order of equal-radius markers}
\begin{lstlisting}[language=C++]
for (std::size_t i = begin + 1; i < end; ++i) {
    if (!(particles[i].r < particles[i - 1].r)) continue;
    SphericalDSMCParticle key = particles[i];
    std::size_t j = i;
    while (j > begin && key.r < particles[j - 1].r) {
        particles[j] = particles[j - 1];
        --j;
    }
    particles[j] = key;
}
\end{lstlisting}
For example, let the radius order after the drift be
\[
 [0.11_A,\;0.15_B,\;0.13_C,\;0.18_D].
\]
When \(C\) is the key, \(B\) moves one place to the right and \(C\) is
inserted before it:
\[
 [0.11_A,\;0.13_C,\;0.15_B,\;0.18_D].
\]
One \emph{inversion} is one out-of-order pair, such as \(B,C\) in this
example.  If the array has \(I\) inversions, insertion sort costs
\(O(N_p+I)\) operations and \(O(1)\) extra marker storage.  After a sufficiently
small physical timestep, most markers remain close to their former radial
rank, \(I\) is small, and insertion sort is almost linear.  Its worst case is
\(O(N_p^2)\), so it is inappropriate after broad radial disorder.

\paragraph{Threaded stable insertion-and-merge path.}
Paper-exact execution must reproduce the stable radial ordering generated by
Frank's insertion sort.  VENUS first counts adjacent inversions in parallel
and counts the 32-separated predicates
\[
 r_i<r_{i-32}\qquad (i\ge32).
\]
This second count is a probe for broad disorder; it is not a stored
measurement of how many ranks a particular marker moved.  With no adjacent
inversion, no sort is needed.  If at most \(N_p/4\) of the 32-separated
predicates are inverted, the paper-exact path applies the literal stable
insertion pass.  Above that threshold, the marker array is divided into
contiguous blocks.  Different threads insertion-sort different blocks, after
which successive stable merge levels combine blocks of width one, two, four,
and so on.  A scratch array holding \(N_p\) complete marker records is used
during the merge.  Choosing a marker from the left block when two radii are
equal preserves the same tie order as the serial stable sort.

For an evolved generic run, VENUS instead counts adjacent inversions.  At
most \(N_p/64\) adjacent inversions select insertion sort; a larger count
selects per-thread stable sorting of contiguous blocks followed by the same
stable merge tree.  A multi-thread initial sort uses radix sorting, while a
one-thread initial sort uses the standard stable general sort.  These
thresholds choose work, not physics: every branch returns the same
nondecreasing stable radius order needed by the force calculation.

\paragraph{Radix sort for a broadly disordered general run.}
A radix is the number of possible digits processed in one pass.  VENUS uses
radix 256, so one digit is one byte and there are 256 buckets.  A
\emph{bucket} is the consecutive output region assigned to one byte value.
All valid radii are finite and nonnegative; for such IEEE-754 double-precision
numbers, the unsigned 64-bit bit pattern has the same order as the numerical
radius.  The code therefore processes the eight bytes of the radius from
least significant to most significant.  Each byte pass is stable, so after
eight passes the complete 64-bit keys are in numerical order.

The parallel byte pass has four steps:
\begin{enumerate}
\item Divide the input into \(T\) consecutive thread chunks.  Thread \(t\)
      counts its byte values in a private histogram \(H_{td}\), where
      \(d=0,\ldots,255\).  A histogram is simply a table of counts.
\item Sum the thread histograms to obtain the total size of each bucket and
      its global starting position \(B_d\).
\item Give thread \(t\) its private starting offset
      \[
       O_{td}=B_d+\sum_{q<t}H_{qd}.
      \]
\item Each thread scans its own input chunk from left to right and writes a
      marker into its current \(O_{td}\), then increments that offset.
\end{enumerate}
The offsets assign disjoint output positions, so scattering requires no
atomic updates or locks.  Earlier input chunks receive earlier positions
inside every bucket, and each thread scans in input order; these two facts
make the parallel pass stable.

As a one-byte example, suppose two threads receive
\[
\begin{split}
 t=0 &: [13_A,2_B,13_C,5_D],\\
 t=1 &: [2_E,9_F,5_G,13_H].
\end{split}
\]
The nonzero local counts are
\[
\begin{array}{c|rrrr}
 &d=2&d=5&d=9&d=13\\ \hline
H_{0d}&1&1&0&2\\
H_{1d}&1&1&1&1
\end{array}
\]
and the global bucket starts are \(B_2=0\), \(B_5=2\), \(B_9=4\), and
\(B_{13}=5\).  Thread 0 therefore starts at positions
\((0,2,4,5)\); thread 1 starts at \((1,3,4,7)\).  Parallel scattering produces
\[
 [2_B,2_E,\;5_D,5_G,\;9_F,\;13_A,13_C,13_H],
\]
including the original order of equal keys.  The actual radius sorter repeats
this operation for all eight bytes:
\codesource{SphDSMC stable radix sorter}
{\path{cpu/include/ev/spherical_dsmc_solver.hpp}}
{Performs the general parallel stable radius ordering used for a broadly
 disordered or explicitly forced sort}
\begin{lstlisting}[language=C++]
for (int pass = 0; pass < 8; ++pass) {
    build_thread_private_histograms(pass);
    build_global_bucket_and_thread_offsets();
    parallel_stable_scatter(source, scratch, pass);
    std::swap(source, scratch);
}
\end{lstlisting}
Radix sorting costs \(O(8N_p)\) marker visits and one \(N_p\)-marker scratch
array.  Because it reads and writes the full 40-byte marker record eight
times, it is usually limited by memory bandwidth and is wasteful when only a
few crossings need repair.  The generic sorter therefore keeps insertion
sort for nearly ordered data.  Radix sorting can be forced for a benchmark
with \code{EV_SPHERICAL_DSMC_RADIX_SORT=1}; it is separate from the
paper-exact post-drift insertion rule.

\subsubsection{Physical global timestep}

The gravity timestep follows the SphDSMC dynamical criterion
\[
 \Delta t_{\rm grav}
 =\left[
 2\eta_{\rm grav}
 \min_i\left(\frac{r_i}{|a_{r,i}|}\right)
 \right]^{1/2},                                            \tag{S6}
\]
where \(\eta_{\rm grav}=\code{bp_dsmc_gravity_eta}\).  This supplies the orbit
accuracy criterion.

For collision shell \(s\), the code finds an upper bound
\({\cal K}_s\) on the physical pair kernel
\(u\Sigma_{ij}(u)\) for every possible marker pair in
that shell.  Collision-sampling literature calls such an upper bound a
\emph{majorant}.  Let
\[
 n_{{\rm partner},s}=\frac{N_s-1}{V_s}
\]
be the finite-shell density of possible partners for one selected marker.
The implemented collision limit is
\[
 \Delta t_{\rm coll}
 =\epsilon_{\rm coll}\min_s
 \begin{cases}
 \displaystyle
 \frac{4}{n_{{\rm partner},s}{\cal K}_s},
     &\text{paper-exact path},\\[1.1ex]
 \displaystyle
 \frac{1}{n_{{\rm partner},s}{\cal K}_s},
     &\text{generic path},
 \end{cases}                                               \tag{S7}
\]
Here
\(\epsilon_{\rm coll}=\code{bp_dsmc_collision_epsilon}\); its default is one.
For equal masses \(m_p\), a constant cross section per unit mass
\(s_\chi\), and the paper majorant, this becomes
\({\cal K}_s=4v_{\max,s}m_ps_\chi\), so the factor four cancels and
\[
 \Delta t_{\rm coll}
 =\epsilon_{\rm coll}
  \min_s\frac{1}{n_{{\rm partner},s}m_ps_\chi v_{\max,s}}.
\]
This is Frank and Tokayer's equation~(12), with the finite-\(N_s\) partner
density retained.  The \(4v_{\max}\) factor belongs to their trial-pair
majorant and must not shorten the paper collision timestep by another factor
of four.  The selected step is
\[
 \Delta t=\min\left(
 \Delta t_{\max},\Delta t_{\rm grav},\Delta t_{\rm coll},
 t_{\rm end}-t\right).                                     \tag{S8}
\]
Scattering is called once per physical global step.  The SphDSMC cadence is
therefore fixed by equations (S7)--(S8).

\subsubsection{Selective Levi--Civita orbit retry}

Very radial trajectories can have a pericentre timescale much shorter than
the global orbital time.  Near pericentre, the terms \(L^2/r^3\) and
\(GM/r^2\) can change rapidly during one ordinary step.

VENUS handles these markers through the following sequence:
\begin{enumerate}
\item Advance every marker with the ordinary kick--drift--kick step.
\item For each marker, compare the orbital energy before and after that step.
\item Accept the ordinary result when the relative energy change satisfies
      \code{vp_orbit_energy_tolerance}.
\item Return a failed marker to its original state and integrate the same
      interval in the regularized coordinate \(\varrho=\sqrt r\).
\item Accept the regularized result after its adaptive integrator converges
      and the energy check passes.
\end{enumerate}

The paper-exact energy criterion is Frank and Tokayer's full equation~(5).
For equal marker mass \(m_p\) and zero-based radial rank \(i\), it is
\[
 {\cal E}_i
 =\frac12v_{r,i}^2+\frac{L_i^2}{2r_i^2}
  -G\left[\frac{M_\bullet}{r_i}
  +m_p\left(\frac{i+1}{r_i}
  +\sum_{j=i+1}^{N_p-1}\frac1{r_j}\right)\right]. \tag{S9}
\]
The \(i+1\) term includes this marker's full spherical self contribution;
\(M_\bullet\) adds the configured central point mass.  The implementation
sums individual masses in both terms when the markers have unequal masses.
The full self-consistent energy of all markers is also tracked as a global
gravity diagnostic.

When the relative defect in (S9) exceeds
\code{vp_orbit_energy_tolerance}, the original marker is retraced using
\[
 \varrho=\sqrt r,\qquad \frac{\dd t}{\dd\tau}=r=\varrho^2. \tag{S10}
\]
For physical orbit equations
\((\dot r,\dot v_r)=F(r,v_r)\), the transformed system is
\[
 \frac{\dd\varrho}{\dd\tau}
   =\frac{\varrho}{2}\dot r,\qquad
 \frac{\dd v_r}{\dd\tau}=r\dot v_r,\qquad
 \frac{\dd t}{\dd\tau}=r.                                 \tag{S11}
\]
The paper-exact path integrates the transformed equations with
kick--drift--kick leapfrog steps in fictitious time, accumulates physical time
with the midpoint value of \(\varrho^2\), and linearly interpolates the final
overshooting step.  A failed equation~(S9) check repeats the original orbit
with the fictitious step halved.  The generic VENUS path retains its adaptive
Dormand--Prince integrator.

The retry is accepted only when the adaptive integrator converges and the
final defect in (S9) meets the requested tolerance.  The history reports the
numbers attempted, accepted,
unresolved, and failed inside the adaptive integrator, as well as the largest
ordinary and final energy defects.  A production result requires zero
unresolved retries.
\code{EV_SPHERICAL_DSMC_FORCE_CONSISTENT_ORBIT_AUDIT=1} is an opt-in
finite-\(N\) diagnostic: it subtracts the half-self term that is present in
the published equation-(5) audit potential but absent from the force used to
advance a marker.  It does not change that force, the published diagnostic,
the total energy, or any collision.  The resolved choice is recorded in
\path{metadata.json}.

\subsubsection{Adaptive radial collision shells and locality}

Collisions should pair markers that occupy the same local region.  The
generic execution path builds radial collision shells after the gravity update:
\begin{enumerate}
\item Start from the radius-sorted marker list.
\item Divide it into blocks containing approximately
      \code{bp_dsmc_shell_particles} markers.
\item Compute each block's radial width, density, evolving orbital velocity
      dispersion, and pair-rate-averaged mean free path.
\item Split a block in half when its width exceeds its mean free path and both
      children contain at least \code{bp_dsmc_min_shell_particles} markers.
\item Apply the same criterion recursively to the children.
\end{enumerate}

Paper-exact mode follows Frank and Tokayer's sequential shell construction.
At the first unused marker, it proposes the next target number of markers,
halves that candidate until its width is local or the population floor is
reached, accepts the shell, and advances to the next unused marker.  The next
candidate starts again at the target population.  The final candidate contains
the remaining consecutive markers, and the faces are their radial extrema,
as in their equation~(6).  Rebuilding either shell construction after gravity
uses the current marker positions.

\codesource{Paper-exact SphDSMC collision shells}
{\path{cpu/include/ev/spherical_dsmc_solver.hpp}}
{Shortened moving-cursor update; \code{describe} abbreviates the source's
 shell-description call, while \code{target} and \code{minimum} are the
 configured populations}
\begin{lstlisting}[language=C++]
while (begin < particles.size()) {
    std::size_t end = std::min(particles.size(), begin + target);
    auto shell = describe(begin, end);
    while (shell.outer_edge - shell.inner_edge > shell.mean_free_path
           && end - begin >= 2 * minimum) {
        end = begin + (end - begin) / 2;
        shell = describe(begin, end);
    }
    shells.push_back(shell);
    begin = end;
}
\end{lstlisting}

For shell edges \(r_{s,-},r_{s,+}\), marker count \(N_s\), and individual
marker masses \(m_i\),
\[
 V_s=\frac{4\pi}{3}(r_{s,+}^3-r_{s,-}^3),\qquad
 \rho_s=\frac{\sum_{i\in s}m_i}{V_s},\qquad
 n_s=\frac{N_s}{V_s},\qquad
 n_{{\rm partner},s}=\frac{N_s-1}{V_s},                    \tag{S12}
\]
\[
 \overline v_{r,s}=\frac{\sum_i m_i v_{r,i}}{\sum_i m_i},
 \qquad
 \sigma_{v,s}^{\,2}=
 \frac{\sum_i m_i[(v_{r,i}-\overline v_{r,s})^2+v_{t,i}^2]}
 {3\sum_i m_i},                                           \tag{S13a}
\]
\[
 \boxed{\displaystyle
 \lambda_s=
 \frac{\langle u_{ij}\rangle_s}
 {n_{{\rm partner},s}
  \langle u_{ij}\Sigma_{ij}(u_{ij})\rangle_s}}.             \tag{S13b}
\]
Here \(v_{r,i}\) and \(v_{t,i}=L_i/r_i\) are radial and tangential speed,
\(\overline v_{r,s}\) is the shell's mass-weighted radial drift,
\(\sigma_{v,s}\) is its one-dimensional orbital velocity dispersion,
\(u_{ij}\) is pair relative speed, and \(\Sigma_{ij}\) is the complete
unequal-mass physical pair area.  The density \(n_s\) estimates all objects,
whereas \(n_{{\rm partner},s}\) excludes the selected object itself; replacing
the latter by \(n_s\) is the large-\(N_s\) continuum approximation.  Angle
brackets average over unordered pairs and their unresolved relative
tangential azimuth.  VENUS uses all pairs when the shell contains at most 512
unordered pairs and otherwise uses 512 deterministic samples.  This
lets the evolving distribution function determine both the velocity scale
and a velocity-dependent cross section.  A single RMS-speed substitution is
not made.

For equal masses \(m_p\) and constant \(s_\chi=\sigma/m\),
\(\Sigma_{ij}=m_p s_\chi\), so equation (S13b) gives
\[
 \lambda_s=\frac{N_s}{N_s-1}\frac{1}{\rho_s s_\chi}.
\]
It approaches Frank and Tokayer's continuum expression
\(1/(\rho_s s_\chi)\) as \(N_s\) grows.  For an equal-mass,
constant-cross-section paper run, the code evaluates this finite-\(N_s\)
form directly.  The published Frank benchmark starts with equal masses, but
\code{EV_SPHERICAL_DSMC_PAPER_EXACT=1} does not reject an unequal-mass
restart.  In that case VENUS uses the unequal pair area and the sampled
velocity-dependent average in equation (S13b).  The shell, cadence,
remainder, and orbit conventions remain paper-exact; unequal masses are a
VENUS extension rather than a claim about the original benchmark.
The recursive decision in the source is
\codesource{SphDSMC adaptive collision-shell builder}
{\path{cpu/include/ev/spherical_dsmc_solver.hpp}}
{Bisects a radial collision region when its width exceeds the sampled physical
 mean free path and both child shells retain the minimum population}
\begin{lstlisting}[language=C++]
if (outer_edge - inner_edge > mean_free_path
    && size >= 2 * minimum) {
    split(begin, middle);
    split(middle, end);
}
\end{lstlisting}

After subdivision, VENUS targets
\[
 \Delta r_s\leq\lambda_s.                                  \tag{S14}
\]
Following Frank's SphDSMC algorithm, if the minimum marker population
prevents (S14), shell refinement stops at that population and the run
continues.  This is a spatial-resolution limit: collisions in such a shell
can transport heat too far.  The maximum \(\Delta r_s/\lambda_s\) and number
of violating shells are therefore written to \path{history.csv} and must be
compared with the resolved core size before late collapse is trusted.
Reduce the minimum population in a convergence sequence until the resolved
density and heat-flow profiles are insensitive to that choice.

\subsubsection{Equal-mass initial resolution at a target central density}
\label{sec:sphdsmc-marker-scaling}

Two length comparisons answer different physical questions.  The numerical
locality requirement is
\[
 \Delta r_{\rm coll}<\lambda_{\rm mfp},
\]
where \(\Delta r_{\rm coll}\) is the width of a collision shell.  Increasing
the marker resolution can improve this inequality.  The core Knudsen
condition is
\[
 r_{\rm core}<\lambda_{\rm mfp}.
\]
It states that the physical core is in the long-mean-free-path regime.
Changing the numerical marker count cannot change this second inequality.
SphDSMC can in principle describe a short-mean-free-path core,
\(\lambda_{\rm mfp}<r_{\rm core}\), but only if its collision shells remain
smaller than \(\lambda_{\rm mfp}\).  A conducting-fluid or kinetic--fluid
description becomes more economical when the core contains many mean free
paths.

The equal-mass resolution scaling follows from the innermost nearly uniform
core.  Let \(M\) be the total represented mass, \(N_p\) the global number of
equal-mass markers, and \(k\) the minimum population of the innermost
collision shell.  One marker represents
\[
 m_p=\frac{M}{N_p}.
\]
If the central shell has density \(\rho\) and characteristic radius
\(\Delta r\), its represented mass gives
\[
 k m_p\simeq \frac{4\pi}{3}\rho\,\Delta r^3,
 \qquad
 \Delta r\simeq
 \left(\frac{3kM}{4\pi N_p\rho}\right)^{1/3}.              \tag{S14a}
\]
Define the physical cross section per unit mass
\[
 s\equiv\frac{\sigma}{m},
 \qquad
 \lambda_{\rm mfp}=\frac{k}{k-1}\frac{1}{\rho s}.
\]
The second equality is the finite-\(k\) mean free path in the innermost
collision shell; \(1/(\rho s)\) is its continuum limit.
The locality ratio is therefore
\[
 \chi\equiv\frac{\Delta r}{\lambda_{\rm mfp}}
 \simeq
 \frac{k-1}{k}\,
 s\left(\frac{3kM}{4\pi N_p}\right)^{1/3}\rho^{2/3}.
                                                               \tag{S14b}
\]
Solving for the required global marker count gives
\[
 \boxed{
 N_{\rm req}\simeq
 \frac{3kM}{4\pi}\,
 \left(\frac{k-1}{k}\right)^3
 s^3\rho_{\rm target}^2\chi_{\rm target}^{-3}}
                                                               \tag{S14c}
\]
and hence the central result
\[
 N_{\rm req}\propto
 \frac{(k-1)^3}{k^2}
 \left(\frac{\sigma}{m}\right)^3\rho_{\rm target}^2
 \ \longrightarrow\
 k\left(\frac{\sigma}{m}\right)^3\rho_{\rm target}^2
 \quad(k\gg1).
                                                               \tag{S14d}
\]
Thus one additional decade in target density costs two decades in equal-mass
markers.  Doubling \(\sigma/m\) costs a factor \(2^3=8\).  In the continuum
approximation, lowering the minimum shell population from \(k=32\) to \(k=8\)
saves a factor of four in marker count and extends the density reach by a
factor of two.  Retaining the finite-\(k\) correction changes those factors
to about \(5.43\) and \(2.33\), respectively; this modest difference is
smaller than the need to demonstrate distribution-function convergence in a
shell containing only eight markers.
Demanding the safer margin \(\chi_{\rm target}=0.5\), rather than marginal
locality at \(\chi_{\rm target}=1\), costs another factor of eight.

\paragraph{Marker-count and memory planning.}
Equation (S14c) should be evaluated in the run's consistent unit system using
the target density, selected minimum shell population, and desired locality
margin.  For an equal-mass velocity-dependent model, the local effective
cross section per mass is
\[
 s_{\rm eff}=
 \frac{\langle u\,s(u)\rangle}{\langle u\rangle}.
\]
Using the largest locally relevant \(s(u)\) gives a conservative upper bound
for initial resource planning, while the evolving pair average in equation
(S13b) supplies the actual locality diagnostic.

The SphDSMC work arrays require approximately 180--210 bytes per marker: the
40-byte particle record is accompanied by energy, potential, rank,
shell-speed, collision-remainder, and sorting workspaces.  A useful estimate
is therefore
\[
 M_{\rm RAM}\simeq190N_p\ {
m bytes}.                    \tag{S14e}
\]
A version-6 restart has a 100-byte header and stores 40 bytes per marker;
paper-exact version 7 adds an 8-byte count and an 8-byte collision-remainder
record per marker.  The marker-count parameter is a 32-bit
integer and cannot exceed \(2.147\times10^9\).  SphDSMC uses one shared-memory
process rather than distributing markers over MPI ranks, so the complete
state and sorting workspace must fit on one node.  Runtime grows at least
linearly with marker count for transport and sorting, and the number of steps
also rises during collapse as the collision and orbital times shorten.

\subsubsection{Reconstructing local three-dimensional velocities}

At a radius, spherical symmetry stores only \(v_r\) and the tangential speed
\(v_t=L/r\); the tangential azimuth is uniformly distributed.  In the
paper-exact path a fresh azimuth \(\phi_i\) is sampled whenever a trial pair
is evaluated:
\[
 \bm v_i=(v_{r,i},v_{t,i}\cos\phi_i,v_{t,i}\sin\phi_i).
                                                               \tag{S15}
\]
In code,
\codesource{SphDSMC local velocity reconstruction}
{\path{cpu/include/ev/spherical_dsmc_solver.hpp}}
{Combines stored radial motion with a sampled tangential azimuth to recover a
 local three-dimensional velocity for a collision trial}
\begin{lstlisting}[language=C++]
velocities[local] = {
    particle.vr,
    vt * std::cos(angle),
    vt * std::sin(angle),
};
\end{lstlisting}
This is a continuous-velocity reconstruction, followed by direct assignment
of the post-collision marker velocity.

At \(r=0\), finite velocity requires \(L=0\) and the distribution is
isotropic. The stored radial departure speed is then the full speed.
Both collision samplers reconstruct an isotropic Cartesian direction there
and retain the outgoing magnitude after scattering. Their stored kinetic
energy therefore equals the conserved Cartesian kinetic energy. The
assignment is
\codesource{SphDSMC stored velocity after an elastic encounter}
{\path{cpu/include/ev/spherical_dsmc_solver.hpp}}
{Preserves the complete outgoing speed at the origin and the radial and
 angular components at positive radius}
\begin{lstlisting}[language=C++]
particle.vr = particle.r == 0.0
    ? std::hypot(velocity.x, std::hypot(velocity.y, velocity.z)) : velocity.x;
particle.L = particle.r * std::hypot(velocity.y, velocity.z);
\end{lstlisting}
The local dispersion gives zero mean projection and one-third of the squared
speed along any chosen axis at the origin. The positive-radius construction
in equation (S15) retains its radial and tangential meanings.

\subsubsection{Sampling the Boltzmann collision rate}

Consider one radial shell with \(N_s\) markers, volume \(V_s\), number
density \(n_s=N_s/V_s\), and possibly unequal marker masses.  The configured
material cross section per unit mass is
\[
 s_\chi(u)\equiv\frac{\sigma(u)}{m}.
\]
This notation applies to the \code{constant}, \code{powerlaw}, and
\code{yukawa_born} laws. The \path{softened_gravity} fit is a
physical-area model with a different unequal-mass scaling, given below.
A value of \(\sigma/m\) fixes the physical pair area only for the
single-species equal-mass problem.  It does not, by itself, derive a mixed
cross section for two different physical species.  The maintained
whole-object extension uses the explicit symmetric closure
\[
 \Sigma_{ij}(u)=\frac{m_i+m_j}{2}s_\chi(u),\qquad
 K_{ij}=u_{ij}\Sigma_{ij}(u_{ij}).                         \tag{S16a}
\]
This law is homogeneous under common refinement and recovers the configured
equal-mass area exactly.  Those two properties do not make it a
first-principles mixed-species interaction.  Production material-SIDM
benchmarks initialize equal masses and elastic scattering leaves them equal.
If \(m_i,m_j\) are distinct physical species, their microphysical
\(\Sigma_{ij}(u)\) must be supplied as additional physics.  If they are merely
unequal numerical sampling weights for one SIDM species, updating both whole
markers cannot simultaneously reproduce both weighted Boltzmann loss rates
and conserve represented pair momentum event by event; a dedicated weighted
DSMC construction is required.  VENUS does not silently claim otherwise.

Under the stated whole-object closure, neither partner is divided into
equal-mass packets.  The paper-exact pair-sampling algorithm is:
\begin{enumerate}
\item Bound every possible \(K_{ij}\) in the shell by \({\cal K}_s\).
      For a constant \(s_\chi\), the code uses
      \({\cal K}_s=4v_{\max,s}m_{\max,s}s_\chi\).
\item Compute the equation-(7) expected trial count and convert it to an
      integer using the persistent equation-(8) remainders.
\item Choose two distinct markers uniformly from the shell.
\item Give both markers fresh random tangential azimuths and compute their
      relative speed \(u_{ij}\).
\item Accept the pair with probability \(K_{ij}/{\cal K}_s\).
\item Scatter each accepted pair elastically in an isotropic outgoing
      direction.
\end{enumerate}

The upper bound is
\[
 {\cal K}_s\geq\max_{i,j}K_{ij},
 \qquad u_{ij}=|\bm v_i-\bm v_j|.                          \tag{S16b}
\]
The equation-(7) expected count is
\[
 \widehat N_{{\rm coll},s}
 =\binom{N_s}{2}\frac{{\cal K}_s\Delta t}{V_s}
 =\frac12(N_s-1)n_s{\cal K}_s\Delta t.                     \tag{S17}
\]
This is the physical unordered-pair count; self-pairs are excluded.  It
uses number density because each selected marker is a complete collision
partner.  Its mass enters through \(\Sigma_{ij}\), not through the pair
proposal probability.

Paper-exact execution uses \(\lfloor\widehat N_{{\rm coll},s}\rfloor\)
trials immediately.  Its fractional part is divided equally among the
shell's persistent particle accumulators \(A_i\).  When
\(\sum_{i\in s}A_i>1\), one extra trial is taken and \(1/N_s\) is subtracted
from every \(A_i\) in the shell.  The accumulators remain attached to
immutable particle IDs as particles change shells and are stored in
paper-exact restart format version 7.  A selected pair is accepted with
probability
\[
 P_{ij}=
 \frac{K_{ij}}{{\cal K}_s}.                                \tag{S18}
\]
For \(m_i=m_j\) and constant \(s_\chi\), this reduces exactly to
\(u_{ij}/(4v_{\max,s})\).
The generic non-paper VENUS path retains the exponential null-collision
clock and velocity-dependent analytic majorants described below.  It is not
used for the Frank benchmark.

Analytic majorants are used for every supported law.  For example,
\[
 s_\chi(u)=A s_0(u^2+u_0^2)^{-n/2}
\]
where \(A=\code{bp_collision_strength}\), \(s_0=\code{bp_sigma0}\),
\(u_0=\code{bp_sigma_u0}\), and \(n=\code{bp_sigma_n}\),
has the interior maximum of \(u s_\chi(u)\) at
\[
 u_*=\frac{u_0}{\sqrt{n-1}}\quad(n>1),                     \tag{S19}
\]
The Born-Yukawa cross section per mass is
\[
 s_{\chi,{\rm Y}}(u)=
 A\frac{4\pi\alpha_\chi^2m_\chi^2}
 {m_\phi^4[1+(m_\chi u/m_\phi)^2]}.                        \tag{S20}
\]
Here \(\alpha_\chi=\code{bp_yukawa_alpha}\),
\(m_\chi=\code{bp_yukawa_mchi}\), and
\(m_\phi=\code{bp_yukawa_mphi}\) are numerical input parameters.  Their
ratio sets a speed scale \(u_{\rm Y}=m_\phi/m_\chi\), and their amplitude
sets the zero-speed cross section per mass
\(s_0=4\pi\alpha_\chi^2m_\chi^2/m_\phi^4\).  Thus (S20) is
\(A s_0/[1+(u/u_{\rm Y})^2]\).  The material collision law interprets
this output in units of cross section per mass; \code{bp_sigma0} is unused
for this selection.  The conversion from a microscopic model is given below.
The product \(u s_{\chi,{\rm Y}}(u)\) has its maximum at
\(u_*=m_\phi/m_\chi\).
The constant, softened-gravity, and endpoint values are treated similarly.
If a sampled kernel exceeds its analytic majorant, the operator raises an
explicit error before applying collisions.

\paragraph{Converting a microscopic Yukawa model to material-SIDM inputs.}
\label{sec:yukawa-material-units}
In natural units \(\hbar=c=1\), let \(m_{\chi,{\rm phys}}\) and
\(m_{\phi,{\rm phys}}\) be the physical particle and mediator masses, and
let \(\alpha_{\rm phys}\) be their dimensionless coupling.  The Born
differential cross section is
\[
 \frac{\dd\sigma_{\rm phys}}{\dd\Omega}
 =\frac{\alpha_{\rm phys}^2m_{\chi,{\rm phys}}^2}
 {\left[m_{\phi,{\rm phys}}^2+
 \tfrac12m_{\chi,{\rm phys}}^2u^2(1-\mu)\right]^2},
 \qquad \mu=\cos\theta.                                  \tag{S20Y1}
\]
Here \(u\) is the relative speed in units of \(c\), and
\(y=(m_{\chi,{\rm phys}}u/m_{\phi,{\rm phys}})^2\) measures the angular
dependence of the scattering.  The zero-speed total cross section is
\(\sigma_0=4\pi\alpha_{\rm phys}^2m_{\chi,{\rm phys}}^2/
m_{\phi,{\rm phys}}^4\).  Integration over solid angle gives the total
cross section; weighting the same integral by \(1-\mu\) gives the
momentum-transfer cross section:
\[
 \sigma_{\rm tot}=\frac{\sigma_0}{1+y},\qquad
 \sigma_{\rm T}=\frac{2\sigma_0}{y^2}
       \left[\ln(1+y)-\frac{y}{1+y}\right].              \tag{S20Y2}
\]
Both approach \(\sigma_0\) as \(y\to0\).  The cross section per mass is
\(s_0=\sigma_0/m_{\chi,{\rm phys}}\), converted to the chosen physical
area and mass units before supplying a simulation input.

Let \(M_{\rm u}\), \(R_{\rm u}\), and \(V_{\rm u}\) be the simulation's
mass, length, and speed units.  The required numerical amplitude and speed are
\[
 s_{0,{\rm code}}=\frac{M_{\rm u}}{R_{\rm u}^2}s_{0,{\rm phys}},
 \qquad
 u_{{\rm Y},{\rm code}}=
 \frac{c}{V_{\rm u}}\frac{m_{\phi,{\rm phys}}}{m_{\chi,{\rm phys}}}.
                                                               \tag{S20Y3}
\]
Choose any positive numerical \(m_{\chi,{\rm code}}\), set
\(m_{\phi,{\rm code}}=u_{{\rm Y},{\rm code}}m_{\chi,{\rm code}}\),
and calibrate the input amplitude by
\[
 \alpha_{{\rm code}}=
 \left(\frac{s_{0,{\rm code}}m_{\phi,{\rm code}}^4}
 {4\pi m_{\chi,{\rm code}}^2}\right)^{1/2}.               \tag{S20Y4}
\]
Using these three numerical values in (S20) gives the intended physical
cross section per mass in native units.  The scattering angle uses the same
speed ratio.  The strength \(A\) then multiplies that calibrated law.
The Eulerian grid collision modules interpret their kernel output as an
area and require their own area normalization.

\codesource{Material-SIDM Yukawa parameter conversion}
{\path{cpu/include/ev/common.hpp};
 \path{cpu/include/ev/finite_angle_boltzmann.hpp}}
{Constructs the input amplitude and speed scale for equations (S20) and
 (S20Y3)--(S20Y4); \code{s0_code} and \code{u_y_code} are the converted
 physical inputs}
\begin{lstlisting}[language=C++]
p.bp_yukawa_mchi = 1.0;
p.bp_yukawa_mphi = u_y_code * p.bp_yukawa_mchi;
p.bp_yukawa_alpha = std::sqrt(
    s0_code * std::pow(p.bp_yukawa_mphi, 4)
    / (4.0 * PI * p.bp_yukawa_mchi * p.bp_yukawa_mchi));
p.bp_collision_strength = 1.0;
\end{lstlisting}

\paragraph{Unequal-mass softened-gravity area.}
The \code{softened_gravity} option is a saturated gravitational
momentum-transfer fit, not a material cross section per unit mass.  Let
\[
 m_{\rm ref}=\code{bp_subhalo_mass},\qquad
 a_{\rm ref}=\code{bp_subhalo_radius},\qquad
 \ln\Lambda=\code{bp_coulomb_log}.
\]
Here \(m_{\rm ref}\) and \(a_{\rm ref}\) are the reference subhalo mass and
radius, and \(\ln\Lambda\) is the Coulomb logarithm that represents the
resolved range of gravitational impact parameters.  Homologous constant-
density radii are
\[
 a_i=a_{\rm ref}\left(\frac{m_i}{m_{\rm ref}}\right)^{1/3},
 \qquad
 a_j=a_{\rm ref}\left(\frac{m_j}{m_{\rm ref}}\right)^{1/3},
 \qquad
 a_{ij}=\frac{a_i+a_j}{2}.                                \tag{S20a}
\]
The reference saturation area is
\[
 \Sigma_{{\rm sat},{\rm ref}}=
 \begin{cases}
 \code{bp_sigma0},&\code{bp_sigma0}>0,\\
 2\pi a_{\rm ref}^2,&\code{bp_sigma0}=0,
 \end{cases}
\]
and the unequal-pair saturation area is
\[
 \Sigma_{{\rm sat},ij}
 =\Sigma_{{\rm sat},{\rm ref}}
  \left(\frac{a_{ij}}{a_{\rm ref}}\right)^2.              \tag{S20b}
\]
Thus, only in this softened-gravity branch, a positive
\code{bp_sigma0} has dimensions of physical area rather than area per mass.
The high-speed gravitational amplitude and area are
\[
 C_{ij}=4\pi G^2m_{\rm ref}^2\ln\Lambda
 \left(\frac{m_i+m_j}{2m_{\rm ref}}\right)^2,\qquad
 \Sigma_{{\rm high},ij}(u)=\frac{C_{ij}}{u^4},            \tag{S20c}
\]
where \(G=\code{vp_gravity}\).  The final physical pair area is the harmonic
interpolation
\[
 \boxed{\displaystyle
 \Sigma_{ij}(u)=
 \frac{A}{
  \Sigma_{{\rm sat},ij}^{-1}
  +\Sigma_{{\rm high},ij}^{-1}(u)}} ,                    \tag{S20d}
\]
with \(A=\code{bp_collision_strength}\).  It approaches
\(A\Sigma_{{\rm sat},ij}\) at low speed and
\(A C_{ij}u^{-4}\) at high speed.  The associated outgoing direction is
sampled isotropically; this maintained option is a scalar transfer fit, not a
resolved gravitational deflection calculation.

\subsubsection{Elastic finite-angle scattering}

For an accepted incoming pair, define
\[
 M=m_1+m_2,\qquad
 \mu_{\rm red}=\frac{m_1m_2}{M},\qquad
 \bm V=\frac{m_1\bm v_1+m_2\bm v_2}{M},\qquad
 \bm u=\bm v_1-\bm v_2,\qquad u=|\bm u|.                   \tag{S21}
\]
After sampling an outgoing unit relative direction \(\hat{\bm n}\),
\[
 \bm v_1'=\bm V+\frac{m_2}{M}u\hat{\bm n},\qquad
 \bm v_2'=\bm V-\frac{m_1}{M}u\hat{\bm n}.                 \tag{S22}
\]
Equation (S22) conserves pair momentum and kinetic energy algebraically:
\[
 m_1\bm v_1'+m_2\bm v_2'
 =m_1\bm v_1+m_2\bm v_2,
\]
\[
 \frac12m_1v_1'^2+\frac12m_2v_2'^2
 =\frac12m_1v_1^2+\frac12m_2v_2^2.                        \tag{S23}
\]
The updated spherical variables are assigned directly:
\codesource{SphDSMC finite-angle collision update}
{\path{cpu/include/ev/spherical_dsmc_solver.hpp};
 \path{cpu/include/ev/spherical_dsmc_collision.hpp}}
{Maps the momentum- and energy-conserving outgoing three-dimensional velocity
 back to persistent radial velocity and angular momentum}
\begin{lstlisting}[language=C++]
particle.vr = velocities[local].x;
particle.L = particle.r * std::hypot(
    velocities[local].y,
    velocities[local].z);
\end{lstlisting}
Equation (S22) supplies the complete direct continuous-velocity update.

For constant, softened power-law, and softened-gravity total cross sections,
the outgoing relative direction is isotropic.  For Born-Yukawa scattering,
\[
 \frac{\dd\sigma}{\dd\Omega}
 =A\frac{\alpha_\chi^2m_\chi^2}
 {\left[
 m_\phi^2+\tfrac12m_\chi^2u^2(1-\mu)
 \right]^2},\qquad \mu=\cos\theta.                         \tag{S24}
\]
The code analytically inverts the cumulative probability in \(\mu\) and
samples the scattering azimuth uniformly.  Integrating (S24) over
\(2\pi\,\dd\mu\) gives (S20), fixing the differential-to-total normalization
used by the sampler.

\subsubsection{Maxwellian equilibrium}

At equilibrium, every collision channel has the same probability as its
reverse channel.  This equality of forward and reverse rates is called
detailed balance.  It follows directly from elastic pair conservation.

For equal masses with common mean velocity \(\bm U\) and one-component
velocity variance \(T\),
\[
 f_M(\bm v)\propto
 \exp\left[-\frac{|\bm v-\bm U|^2}{2T}\right],
\]
pair momentum and kinetic-energy conservation imply
\[
 f_M(\bm v_1)f_M(\bm v_2)
 =f_M(\bm v_1')f_M(\bm v_2').                              \tag{S25}
\]
Reversing the outgoing relative direction in (S22) recovers the incoming
pair, and equations (S17)--(S18) sample the physical factor
\(u\,\dd\sigma/\dd\Omega\).  Detailed balance therefore makes the
equal-mass Maxwellian an exact stationary distribution of the underlying
continuous collision process.  For unequal physical object masses, the
equilibrium product instead uses the mass-weighted exponent derived after
(PIC35M), with a common thermodynamic temperature.
The finite-marker realization fluctuates statistically, but its ensemble
residual decreases with \(N_p\).

For a quantitative equilibrium reference, consider \(N\) equal-mass markers
of mass \(m_p\) in volume \(V\).  Define
\[
 n=\frac{N}{V},\qquad
 n_{\rm partner}=\frac{N-1}{V},\qquad
 \rho=nm_p,\qquad
 \Sigma_0=m_ps_0 ,
\]
where \(s_0=\sigma/m\) is the constant cross section per unit mass and
\(\Sigma_0\) is the physical pair area.  The analytic accepted-event rate per
volume is
\[
 {\cal R}_{\rm event}
 =\frac12 n\,n_{\rm partner}\Sigma_0\langle u\rangle
 =\frac12\frac{\rho^2s_0}{m_p}
  \left(1-\frac1N\right)\langle u\rangle,\qquad
 \left\langle u\right\rangle=4\sqrt{\frac{T}{\pi}}.         \tag{S26}
\]
Here \(T\) is the variance of each Cartesian velocity component in equation
(S25).  The factor \(1-1/N\) is the exact finite-ensemble exclusion of
self-pairs.  The sampled rate should approach equation (S26) within its
Monte--Carlo uncertainty, and independent realizations at successively larger
\(N_p\) should show the expected \(N_p^{-1/2}\) decrease of the velocity-
histogram fluctuation.  A fixed numerical value from one random ensemble is
not a physical calibration.

\subsubsection{Conservation ledgers}

Each persistent marker carries its own fixed represented mass throughout the
elastic SphDSMC run.  Every operator step writes explicit before/after
mass-weighted budgets to
\path{operator_mass_budget.csv}.  The fields include
\begin{itemize}[itemsep=0.25em]
\item transport mass and full energy before and after gravity;

\item collision mass before and after scattering;

\item collision radial momentum before and after scattering;

\item collision kinetic energy before and after scattering;

\item maximum pair-level momentum and kinetic-energy residuals.
\end{itemize}
The stored global spherical observables are radial momentum and full pair
kinetic energy.  The tangential azimuth in (S15) supplies an auxiliary local
realization.

The collision step must preserve marker mass exactly, and its mass-weighted
pair momentum and kinetic energy must close to floating-point roundoff.
Transport energy is an integration diagnostic rather than an exact
kick--drift--kick invariant; it must converge as the gravity timestep is
reduced.  A physically interpreted run additionally requires
\(\max_s\Delta r_s/\lambda_s<1\) and zero unresolved LC retries.  These
requirements belong to every run and are not replaced by one stored benchmark
number.

\subsubsection{Initial-condition provenance and acceptance checks}

The consuming constant-law task is
\code{venus_cpu_sidm_sphdsmc_constant}, selected by
\path{cpu/src/apps/sidm/sphdsmc_constant.cpp}.  In the grid-backed route,
\code{read_initial_data_rwl} in \path{cpu/include/ev/initial_data.hpp} imports
an \code{EVICRWL1} object such as
\path{cpu/examples/initial_data/frank_core_eq2_r1e3_model1e6_originlinear_Nr64_Nw32_NL40.bin}.
\code{EVICRWL1} is the bounded magic identifier at the start of a versioned
binary carrier containing the \((r,w,L)\) coordinate grids and the sampled
distribution array \(f(r,w,L)\); it is a file format, not an evolution
variable.  Then \code{spherical_dsmc_sample_initial_particles} in
\path{cpu/include/ev/spherical_dsmc_solver.hpp} samples the in-memory
\code{std::vector<SphericalDSMCParticle>} that becomes the dynamical state.
This C++ object is a contiguous variable-length list of the marker records
defined above.

The solver writes \code{VENUSSPHDSMC1} restart objects beneath the configured
\path{restart_dir}.  \code{VENUSSPHDSMC1} is the bounded magic identifier for
the versioned SphDSMC checkpoint format; the checkpoint contains evolution
time and step metadata together with the persistent marker state and the
version-specific deterministic collision remainder data.
A direct marker route is provided by
\path{cpu/tools/generate_frank_core_spherical_dsmc_restart.py}.  That generator
writes the same \code{VENUSSPHDSMC1} object consumed through
\code{--restart-from}, bypassing carrier-grid sampling.

Whichever route is used, validate the file identifier and version, finite
positive individual masses, unique identifiers, total represented mass,
radial support, the virial ratio \(Q_{\rm vir}=2K/|W|\), energy diagnostics,
initial equilibrium drift, collision-shell occupancy, and mean-free-path
locality.  Repeat the calculation under marker, seed, and timestep refinement;
agreement of one generated state is not a convergence test.

\subsubsection{Restart format and determinism}

SphDSMC restart formats 6 and 7 store each marker's radius, radial velocity,
angular momentum, ID, and individual mass as fixed-width native binary values.
Their header also stores the next collision timestep and a fingerprint of
the collision settings that generated it.  Version 7 includes the per-particle
integer+remainder state used by paper-exact collisions.

The next collision limit is part of the evolving numerical state.  After a
collision stage, the solver computes this limit from that stage's shell
geometry and retained pair-kernel majorants.  A fresh shell reconstruction
from the scattered particles can give a different limit, particularly when
paper-exact scattering changes the largest speed.  Restoring the stored
limit therefore preserves the next accepted timestep.  The implementation
selects the stored value before it enters the timestep loop:
\codesource{SphDSMC collision-clock restoration}
{\path{cpu/include/ev/spherical_dsmc_solver.hpp}}
{Restores the carried collision timestep and reconstructs an initial limit
 for a fresh ensemble or a legacy checkpoint}
\begin{lstlisting}[language=C++]
if (restored_collision_dt.has_value()) {
    next_collision_dt = restored_collision_dt.value();
} else {
    const std::vector<SphericalDSMCShell> initial_shells =
        spherical_dsmc_build_shells(
            particles, particle_mass, state.global_p,
            timestep_shell_workspace);
    next_collision_dt = spherical_dsmc_collision_timestep(
        initial_shells, state.global_p);
    if (restart_input_version != 0) {
        std::cout << "spherical_dsmc legacy restart version="
                  << restart_input_version
                  << " has no persisted collision clock; "
                     "reconstructing the first collision limit\n";
    }
}
\end{lstlisting}
The fingerprint checks the cross-section law and its parameters, angular
kernel, collision strength, shell target and floor, collision timestep
fraction, random seed, and collision mode.  It also includes the gravitational
constant used by the softened-gravity pair law.  Changing these settings
requires a separately initialized physical case; an exact continuation
rejects a mismatched clock.  An infinite positive limit represents an
inactive collision constraint.

Legacy versions 2--5 remain readable.  Versions 2 and 3 copy their common
header mass onto every restored marker; versions 4 and 5 already store
individual masses.  These legacy files initialize the first collision limit
from the restored ensemble, and metadata records this reconstruction.  Their
subsequent checkpoints use the complete current state.  This import preserves
the stored physical ensemble and the available fractional remainders; exact
reproduction of the missing historical collision clock requires a current
checkpoint.

The reader validates the bounded magic identifier, supported version,
nonnegative finite time and safely representable step, initial energy and
counters, positive finite masses, finite coordinates, nonnegative \(L\), and
the uniqueness and range of all stable marker IDs.  Before allocating the
particle vector, it proves that the count fits the version-specific payload.
The complete payload must end after its declared particle and remainder
records, and a stored collision limit must be positive and non-NaN.
Because the writer does not canonicalize byte order or floating-point
representation, move a checkpoint between machines only when their binary
representations are compatible; the CSV diagnostics are the portable
exchange format.

The checkpoint is first written to a temporary file and atomically renamed.
The random streams are deterministic functions of the base seed, global step,
and shell index.  With unchanged physical and numerical settings,
continuation at a natural step boundary preserves the stored stochastic
state, collision clock, and marker ordering.  The actual-solver regression
compares six uninterrupted steps with a three-step run followed by three
resumed steps in both collision modes, including accepted collisions and the
complete final checkpoint.

\subsubsection{Shared-memory organization and computational scaling}

The module requires one process:
\codesource{SphDSMC process-layout contract}
{\path{cpu/Makefile}; \path{cpu/include/ev/parallel.hpp};
 \path{cpu/include/ev/spherical_dsmc_solver.hpp}}
{Documents the enforced one-process execution model and assigns parallel work
 to portable C++ threads or OpenMP threads inside that process}
\begin{lstlisting}
#SBATCH --ntasks=1
#SBATCH --cpus-per-task=<OpenMP threads>
export OMP_NUM_THREADS=$SLURM_CPUS_PER_TASK
\end{lstlisting}
The unsuffixed executable uses the portable C++ \code{std::thread} backend;
the \code{_openmp} executable uses OpenMP.  Both run as one process.  Gravity
loops are parallel over markers.  Collision shells are disjoint and
dynamically scheduled over the selected thread backend, giving each marker
one thread owner during a shell update.

The stochastic work is proportional to the number of physically sampled NTC
trials.  Each accepted event realizes one continuous outgoing angle and
directly updates its two markers.  Marker drifts, force evaluations, profile
deposition, inversion counting, radix histograms, and independent collision
shells are parallel.  The stable radius merge, enclosed-mass prefix work,
suffix potential, and memory movement contain serial or memory-bandwidth
limits, so doubling the thread count need not halve the elapsed time.  For
\(N_p\) markers, state and sorting scratch require \(O(N_p)\) memory, while
one ordinary step costs \(O(N_p+N_{\rm trial})\) after sorting, where
\(N_{\rm trial}\) is the NTC trial count.  Measure a short fixed-work
strong-scaling run at 1, 2, 4, 8, 16, and 32 threads on the target machine and
choose the smallest allocation beyond which the elapsed time changes little.

\subsubsection{Configuration block}

A constant-cross-section SphDSMC run begins with the task selectors, physical
controls, initial-data source, time limits, input/output (I/O) paths, and
thread count below:
\codesource{\code{venus_cpu_sidm_sphdsmc_constant}}
{\path{cpu/src/apps/sidm/sphdsmc_constant.cpp};
 \path{cpu/include/ev/config.hpp}}
{Parses the explicit constant-law SphDSMC configuration and rejects task-
 incompatible physics selections}
\begin{lstlisting}
[physics]
physics_model = boltzmann_poisson
vp_universe = static
vp_gravity = 1
vp_poisson_source = total_density
vp_orbit_regularization = auto
vp_orbit_energy_tolerance = 0.01
vp_orbit_lc_relative_tolerance = 1e-9
vp_orbit_lc_absolute_tolerance = 1e-12
vp_orbit_lc_max_steps = 2048
vp_orbit_lc_max_refinements = 20

[collisions]
bp_collision_operator = spherical_dsmc
bp_cross_section = constant
bp_boltzmann_kernel = isotropic
bp_collision_strength = 1
bp_sigma0 = 1
bp_dsmc_particles = 1000000
bp_dsmc_total_mass = 1.1252580933331495
bp_dsmc_shell_particles = 4096
bp_dsmc_min_shell_particles = 32
bp_dsmc_collision_epsilon = 1
bp_dsmc_gravity_eta = 0.01
bp_dsmc_seed = 5489

[initial_data]
initial_data = binary_rwl
initial_data_file = examples/initial_data/frank_core_eq2_r1e3_model1e6_originlinear_Nr64_Nw32_NL40.bin
custom_ic_normalize_mass = false

[time_integration]
integrator = semi_lagrangian_kick_drift_kick
t_end = 0.05
dt_max = 0.01
max_steps = 1000000

[output]
output_dir = output/spherical_dsmc_example
restart_dir = output/spherical_dsmc_example/restarts
save_every = 1
phase_every = 0
restart_every = 5

[parallelism]
threads = 8
\end{lstlisting}
Because this example selects \code{bp_cross_section = constant},
\code{bp_sigma0} is interpreted as cross section per unit mass in the chosen
code units, so the continuum local rate is \(\rho u(\sigma/m)\).  The value
one above is a runnable dimensionless example.  For
\code{bp_cross_section = softened_gravity}, however, a positive
\code{bp_sigma0} is the physical reference saturation area in equation
(S20b), not a cross section per mass.  A literature benchmark must convert
its quoted cross section, density, velocity, length, and time into one
consistent unit system before replacing these values.

Build the constant-cross-section task frontend and run the parameter file from
the CPU source directory:
\codesource{\code{venus_cpu_sidm_sphdsmc_constant}}
{\path{cpu/src/apps/sidm/sphdsmc_constant.cpp};
 \path{cpu/Makefile}}
{Builds and runs the fail-closed constant-isotropic SphDSMC task frontend}
\begin{lstlisting}[basicstyle=\ttfamily\scriptsize]
cd VENUS/cpu
make venus_cpu_sidm_sphdsmc_constant
./venus_cpu_sidm_sphdsmc_constant --config params.ini
\end{lstlisting}
The executable's \code{--help} interface lists the collision, orbit, shell
and restart controls. A fresh calculation reads its initial distribution;
a continuation reads the complete persistent-marker restart. The physical
SIDM demonstrations and saved profiles are described in \path{analysis/README.md}.

\subsubsection{Convergence requirements}

A gravothermal-collapse calculation requires a joint refinement of the
marker, orbit, collision, and output scales:
\begin{enumerate}
\item increase the marker count by factors of two and compare density,
velocity dispersion, heat flux, collision rate, and core radius;
\item halve \(\epsilon_{\rm coll}\), the gravity parameter, and the target
shell population independently, then jointly;
\item require pair momentum and kinetic energy to close to floating-point
accuracy, transport energy error to decrease with timestep, all orbit retries
to resolve, and \(\Delta r_s/\lambda_s<1\) in every collision shell;
\item compare constant and velocity-dependent kernels with equation (S26),
detailed balance, the analytic majorants, and the Yukawa differential-to-total
normalization where applicable;
\item compare uninterrupted and restarted evolution at a natural step
boundary and repeat stochastic profile comparisons with independent seeds;
\item compare with the selected literature density profile through the
      full quoted dimensionless time;
\item follow a constant-cross-section Plummer core-formation and core-collapse
      sequence with the same unit conversion and boundary convention.
\end{enumerate}
